% cover-tikz.tex — LeetCUDA 书封面（CuTe layout 风格, 浅色 A4 满幅）
% 构建: xelatex cover-tikz.tex (standalone 输出 A4 PDF)
% 主网格 L = (8,16):(1,8), cell(i0,i1) = i0 + 8*i1 (gen_grid.py assert 双射 0..127);
% 配色 = colfax quotient ramp: hue 按 2 列 tile 组 (colorx1..8), intensity 行渐变 8..64
\documentclass[tikz,border=0pt]{standalone}
\usetikzlibrary{arrows.meta,calc,decorations.pathmorphing}
\usepackage{xeCJK}
\usepackage{amsmath}
\setmainfont{Humor Sans}
\setsansfont{Humor Sans}
\setmonofont{Comic Neue}
\setCJKmainfont{LXGW WenKai}
\setCJKsansfont{LXGW WenKai}

\tikzset{
  sketch/.style={decorate, decoration={random steps, segment length=2.5mm, amplitude=0.3mm}},
  sketchfine/.style={decorate, decoration={random steps, segment length=8mm, amplitude=0.15mm}}}

\definecolor{coverBg}{RGB}{248,249,251}
\definecolor{coverInk}{RGB}{31,35,40}
\definecolor{coverMuted}{RGB}{118,126,138}
\definecolor{accentBlue}{RGB}{33,113,181}
% 8 色 jewel-tone (降饱和版, 浅底雅致; 与网格/色条共用)
\definecolor{colorx1}{RGB}{213,84,76}
\definecolor{colorx2}{RGB}{232,163,61}
\definecolor{colorx3}{RGB}{155,191,90}
\definecolor{colorx4}{RGB}{76,182,149}
\definecolor{colorx5}{RGB}{69,168,200}
\definecolor{colorx6}{RGB}{62,124,192}
\definecolor{colorx7}{RGB}{126,99,184}
\definecolor{colorx8}{RGB}{194,88,158}

\begin{document}
\begin{tikzpicture}
% ===== A4 背景 (cm), 极淡冷灰渐变 =====
\shade[top color=coverBg, bottom color=coverBg!94!accentBlue] (0,0) rectangle (21,29.7);

% ===== 背景装饰: 散落 size:stride 元组纹理 (浅灰) =====
\node[rotate=12, anchor=center, text=coverMuted!35, font=\ttfamily\footnotesize] at (2.3,24.3) {(4,8):(8,1)};
\node[rotate=-12, anchor=center, text=coverMuted!35, font=\ttfamily\footnotesize] at (18.7,24.3) {(128,128):(1,128)};
\node[rotate=7, anchor=center, text=coverMuted!30, font=\ttfamily\footnotesize] at (2.3,17.6) {((2,2),(2,2))};
\node[rotate=-7, anchor=center, text=coverMuted!30, font=\ttfamily\footnotesize] at (18.7,17.6) {(3):(1)};
\node[rotate=-6, anchor=center, text=coverMuted!28, font=\ttfamily\footnotesize] at (2.3,7.0) {(16,8):(8,1)};
\node[rotate=6, anchor=center, text=coverMuted!28, font=\ttfamily\footnotesize] at (18.7,7.0) {(8,16):(1,8)};

% ===== 顶部小标签 =====
\node[text=coverMuted, font=\ttfamily\bfseries\small] at (10.5,26.9)
  {CUTE\; LAYOUT $\cdot$ TENSOR\; CORE $\cdot$ FLASH\; ATTENTION};

% ===== 标题块 =====
\node[text=coverInk, font=\fontsize{44}{48}\selectfont\bfseries] at (10.5,23.9) {LeetCUDA};
\node[text=coverInk, font=\fontsize{25}{30}\selectfont\bfseries] at (10.5,21.3)
  {\CJKfamily{zhsans}\sffamily CUDA Kernel 优化之路};
\node[text=coverMuted, font=\fontsize{13}{16}\selectfont] at (10.5,19.5)
  {从编程模型到 CuTe 与 FlashAttention};

% colorx 八段细色条 (与网格 8 hue 呼应, 极细低调)
\foreach \i in {0,...,7} {%
  \fill[colorx\number\numexpr\i+1\relax!85] (3.3+\i*1.8,18.55) rectangle +(1.8,0.06);}

% ===== 主视觉: CuTe layout 网格 (8x16) =====
% L = (8,16):(1,8); hue = 2 列 tile 组, intensity = 行渐变 (quotient ramp)
\begin{scope}[shift={(3.75,15.5)}, x={(0cm,-0.9cm)}, y={(0.9cm,0cm)},
  every node/.style={minimum size=0.9cm, outer sep=0pt, text=coverInk!65,
    font=\fontsize{7.5}{8.5}\selectfont}]
  \node[fill=colorx1!8] at (0,0) {0};
\node[fill=colorx1!8] at (0,1) {8};
\node[fill=colorx2!8] at (0,2) {16};
\node[fill=colorx2!8] at (0,3) {24};
\node[fill=colorx3!8] at (0,4) {32};
\node[fill=colorx3!8] at (0,5) {40};
\node[fill=colorx4!8] at (0,6) {48};
\node[fill=colorx4!8] at (0,7) {56};
\node[fill=colorx5!8] at (0,8) {64};
\node[fill=colorx5!8] at (0,9) {72};
\node[fill=colorx6!8] at (0,10) {80};
\node[fill=colorx6!8] at (0,11) {88};
\node[fill=colorx7!8] at (0,12) {96};
\node[fill=colorx7!8] at (0,13) {104};
\node[fill=colorx8!8] at (0,14) {112};
\node[fill=colorx8!8] at (0,15) {120};
\node[fill=colorx1!16] at (1,0) {1};
\node[fill=colorx1!16] at (1,1) {9};
\node[fill=colorx2!16] at (1,2) {17};
\node[fill=colorx2!16] at (1,3) {25};
\node[fill=colorx3!16] at (1,4) {33};
\node[fill=colorx3!16] at (1,5) {41};
\node[fill=colorx4!16] at (1,6) {49};
\node[fill=colorx4!16] at (1,7) {57};
\node[fill=colorx5!16] at (1,8) {65};
\node[fill=colorx5!16] at (1,9) {73};
\node[fill=colorx6!16] at (1,10) {81};
\node[fill=colorx6!16] at (1,11) {89};
\node[fill=colorx7!16] at (1,12) {97};
\node[fill=colorx7!16] at (1,13) {105};
\node[fill=colorx8!16] at (1,14) {113};
\node[fill=colorx8!16] at (1,15) {121};
\node[fill=colorx1!24] at (2,0) {2};
\node[fill=colorx1!24] at (2,1) {10};
\node[fill=colorx2!24] at (2,2) {18};
\node[fill=colorx2!24] at (2,3) {26};
\node[fill=colorx3!24] at (2,4) {34};
\node[fill=colorx3!24] at (2,5) {42};
\node[fill=colorx4!24] at (2,6) {50};
\node[fill=colorx4!24] at (2,7) {58};
\node[fill=colorx5!24] at (2,8) {66};
\node[fill=colorx5!24] at (2,9) {74};
\node[fill=colorx6!24] at (2,10) {82};
\node[fill=colorx6!24] at (2,11) {90};
\node[fill=colorx7!24] at (2,12) {98};
\node[fill=colorx7!24] at (2,13) {106};
\node[fill=colorx8!24] at (2,14) {114};
\node[fill=colorx8!24] at (2,15) {122};
\node[fill=colorx1!32] at (3,0) {3};
\node[fill=colorx1!32] at (3,1) {11};
\node[fill=colorx2!32] at (3,2) {19};
\node[fill=colorx2!32] at (3,3) {27};
\node[fill=colorx3!32] at (3,4) {35};
\node[fill=colorx3!32] at (3,5) {43};
\node[fill=colorx4!32] at (3,6) {51};
\node[fill=colorx4!32] at (3,7) {59};
\node[fill=colorx5!32] at (3,8) {67};
\node[fill=colorx5!32] at (3,9) {75};
\node[fill=colorx6!32] at (3,10) {83};
\node[fill=colorx6!32] at (3,11) {91};
\node[fill=colorx7!32] at (3,12) {99};
\node[fill=colorx7!32] at (3,13) {107};
\node[fill=colorx8!32] at (3,14) {115};
\node[fill=colorx8!32] at (3,15) {123};
\node[fill=colorx1!40] at (4,0) {4};
\node[fill=colorx1!40] at (4,1) {12};
\node[fill=colorx2!40] at (4,2) {20};
\node[fill=colorx2!40] at (4,3) {28};
\node[fill=colorx3!40] at (4,4) {36};
\node[fill=colorx3!40] at (4,5) {44};
\node[fill=colorx4!40] at (4,6) {52};
\node[fill=colorx4!40] at (4,7) {60};
\node[fill=colorx5!40] at (4,8) {68};
\node[fill=colorx5!40] at (4,9) {76};
\node[fill=colorx6!40] at (4,10) {84};
\node[fill=colorx6!40] at (4,11) {92};
\node[fill=colorx7!40] at (4,12) {100};
\node[fill=colorx7!40] at (4,13) {108};
\node[fill=colorx8!40] at (4,14) {116};
\node[fill=colorx8!40] at (4,15) {124};
\node[fill=colorx1!48] at (5,0) {5};
\node[fill=colorx1!48] at (5,1) {13};
\node[fill=colorx2!48] at (5,2) {21};
\node[fill=colorx2!48] at (5,3) {29};
\node[fill=colorx3!48] at (5,4) {37};
\node[fill=colorx3!48] at (5,5) {45};
\node[fill=colorx4!48] at (5,6) {53};
\node[fill=colorx4!48] at (5,7) {61};
\node[fill=colorx5!48] at (5,8) {69};
\node[fill=colorx5!48] at (5,9) {77};
\node[fill=colorx6!48] at (5,10) {85};
\node[fill=colorx6!48] at (5,11) {93};
\node[fill=colorx7!48] at (5,12) {101};
\node[fill=colorx7!48] at (5,13) {109};
\node[fill=colorx8!48] at (5,14) {117};
\node[fill=colorx8!48] at (5,15) {125};
\node[fill=colorx1!56] at (6,0) {6};
\node[fill=colorx1!56] at (6,1) {14};
\node[fill=colorx2!56] at (6,2) {22};
\node[fill=colorx2!56] at (6,3) {30};
\node[fill=colorx3!56] at (6,4) {38};
\node[fill=colorx3!56] at (6,5) {46};
\node[fill=colorx4!56] at (6,6) {54};
\node[fill=colorx4!56] at (6,7) {62};
\node[fill=colorx5!56] at (6,8) {70};
\node[fill=colorx5!56] at (6,9) {78};
\node[fill=colorx6!56] at (6,10) {86};
\node[fill=colorx6!56] at (6,11) {94};
\node[fill=colorx7!56] at (6,12) {102};
\node[fill=colorx7!56] at (6,13) {110};
\node[fill=colorx8!56] at (6,14) {118};
\node[fill=colorx8!56] at (6,15) {126};
\node[fill=colorx1!64] at (7,0) {7};
\node[fill=colorx1!64] at (7,1) {15};
\node[fill=colorx2!64] at (7,2) {23};
\node[fill=colorx2!64] at (7,3) {31};
\node[fill=colorx3!64] at (7,4) {39};
\node[fill=colorx3!64] at (7,5) {47};
\node[fill=colorx4!64] at (7,6) {55};
\node[fill=colorx4!64] at (7,7) {63};
\node[fill=colorx5!64] at (7,8) {71};
\node[fill=colorx5!64] at (7,9) {79};
\node[fill=colorx6!64] at (7,10) {87};
\node[fill=colorx6!64] at (7,11) {95};
\node[fill=colorx7!64] at (7,12) {103};
\node[fill=colorx7!64] at (7,13) {111};
\node[fill=colorx8!64] at (7,14) {119};
\node[fill=colorx8!64] at (7,15) {127};
  % 手绘风: 只画内部网格线; 外边界由墨线框覆盖 (避免白线与墨框双层错位)
  \foreach \a in {0.5,1.5,...,6.5} \draw[white, line width=0.7pt, sketchfine] (\a,-0.5) -- (\a,15.5);
  \foreach \b in {0.5,1.5,...,14.5} \draw[white, line width=0.7pt, sketchfine] (-0.5,\b) -- (7.5,\b);
  \draw[coverInk!35, line width=0.7pt, sketch] (-0.5,-0.5) rectangle (7.5,15.5);
\end{scope}

% 高亮 tile B = (2,2):(1,8) (右上角 2x2 tile, accent 粗框; scope 内 i0 行 0-1, i1 列 14-15)
\draw[accentBlue, line width=1.5pt, rounded corners=1.5pt, sketch,
  shift={(3.75,15.5)}, x={(0cm,-0.9cm)}, y={(0.9cm,0cm)}]
  (-0.5,13.5) rectangle (1.5,15.5);
\node[anchor=south, text=accentBlue, font=\ttfamily\bfseries\footnotesize]
  at (16.8,16.85) {B = (2,2):(1,8)};
\draw[-{Stealth[length=1.8mm]}, accentBlue, line width=0.8pt, sketch] (16.8,16.7) -- (16.8,16.0);

% layout 方程 (网格正下方, 居中)
\node[text=coverInk, font=\fontsize{13}{16}\selectfont] at (10.5,7.7)
  {L = (8,16):(1,8)};
% mode 标注: mode0 (stride 1) 向下, mode1 (stride 8) 向右
\node[rotate=90, anchor=center, text=coverMuted, font=\ttfamily\footnotesize] at (2.95,12.35) {mode\,0\,:\,(1)\;$\downarrow$};
\node[text=coverMuted, font=\ttfamily\footnotesize] at (10.5,16.35) {$\rightarrow$\;mode\,1\,:\,(8)};

% TiledMMA 代码块 (单行; 卡片与网格同宽对齐, 编辑器窗口风格)
\draw[fill=white!30!coverBg!95!accentBlue, draw=coverInk!35, line width=0.6pt, rounded corners=2.5pt, sketch]
  (3.3,6.0) rectangle (17.7,7.0);
\foreach \i/\c in {0/colorx1,1/colorx2,2/colorx4}
  \fill[\c!80] (3.62+\i*0.2,6.83) circle (0.052);
\node[anchor=center, text=coverMuted, font=\ttfamily\small] at (10.5,6.36)
  {auto tiled\_mma = {\color{accentBlue}make\_tiled\_mma}(SM80\_16x8x16\_F32F16F16F32\_TN\{\}, Layout<Shape<\_2,\_2,\_1>>\{\});};

% ===== 底部声明 (细线框) =====
\node[draw=accentBlue!45, sketch, rounded corners=2pt, inner xsep=10pt, inner ysep=8pt,
  text=coverInk, text width=13.7cm, align=center, font=\small] at (10.5,4.75)
  {\CJKfamily{zhsans}\bfseries 郑重声明\\
   \mdseries\CJKfamily{zhsong} 本书仅用于开源学习与技术交流，\\未经作者本人同意，不得用于任何商业用途。};
\node[text=coverMuted, font=\small] at (10.5,2.85)
  {LeetCUDA Project (Own by DefTruth) $\cdot$ 2026 年 9 月};
\end{tikzpicture}
\end{document}
